\par\noindent\textbf{(1)} 用正反行波振幅(A⁺,A⁻)，无损互易单片可写M=[α,β;β*,α*]，其中|α|²−|β|²=1，因而detM=1。若单片透射概率T₁=1−r，可取|α|²=1/T₁、|β|²=r/T₁；转移矩阵乘法而非透射幅简单相乘才是多片传播的正确描述。
\par\noindent\textbf{(2)} 两散射体间有传播相位P=diag(e\textasciicircum{}\{iφ\},e\textasciicircum{}\{-iφ\})，总矩阵M₂PM₁；分母的多次反射项随φ干涉，故总反射功率不能写成r₁+r₂。对φ均匀随机时对lnT而非T直接平均可消去干涉相位项，得到可加的Lyapunov指数。
\par\noindent\textbf{(3)} 在弱散射并且独立相位近似中，每片平均对数透射 ⟨lnT₁⟩=ln(1−r)≈−r，N片给 ⟨lnT\_N⟩≈Nln(1−r)≈−Nr。令长度L=Na，以标准定义T\_typ=exp⟨lnT⟩=exp(−2L/ξ)，得 ξ≈2a/r。因r小需同时有足够大N才出现强局域。
\par\noindent\textbf{(4)} 在一维随机系统，lnT可出现广分布，故 ⟨T⟩通常不等于exp⟨lnT⟩；长样品的平均透射可能由极罕见高透射实现主导。ξ由典型对数衰减定义，并非每个样本都有严格相同衰减长度；相关随机相位或吸收会破坏上述简单ξ。
